\scanpage{26}
\mspage{7}

On va maintenant appliquer la prop. 1$^{*}$, pour calculer $\pi_1(X,e)$, en termes de $\pi_1(X_0,e)=G_0$, et des relations explicitées par ailleurs. Je veux \deleted{\unsure{prouver}} \inserted{montrer} que les images de ces relations dans $\mathbb Z$ par $\delta_0$ sont nulles, donc que $\pi_1(X_0,e)$ admet un groupe quotient isomorphe à $\mathbb Z$, donc n'est pas $\sim\{1\}$, donc que $X_0$ n'est pas $1$-connexe.
% slanted
\margin{Si on applique le NB page 2, \deleted{\illegible{1}} les développements jusqu'à la page 12 sont inutiles !}

Je reprends donc le diagramme
\diagram{\begin{tikzcd}[ampersand replacement=\&] \& x\arrow[r,"f"]\&y\arrow[r,"g"]\&z\\ \text{(*)}\&\&e\arrow[ul,"u"]\arrow[u,"v"]\arrow[ur,"w"']\&{,}\end{tikzcd}}

en posant
\[
\left\{\begin{aligned}z&=a\rho,\quad\text{donc}\\ y&=a\rho g\\ x&=a\rho gf\end{aligned}\right.
\]
\[
\left\{\begin{aligned}e&=xu=a\rho gfu\\ e&=yv=a\rho gv\\ e&=zw=a\rho w\end{aligned}\right.
\]
et comme $e=a\psi$, les trois relations précédentes s'écrivent
\[
\tag*{(**)}\boxed{\rho gfu\sim\rho gv\sim\rho w\sim\psi}
\]

La donnée d'un diagramme \unsure{cas|($*$)} dans $X$, avec $z=a\rho$ (pour $\rho\in M$ fixé) équivaut à la donnée des six éléments

% archive sheet number: 25
